\subsection{附属于挠模的除子}

沿用第 2、3、4 节的记号和假设。回忆，D(A)（或简称 D）表示 A 的除子群，用加法记号书写；我们知道（§ 1，第 3 节，定理 2），D 是由 P 的元素生成的自由 Z-模。

设 T 是有限生成的挠 A-模。对于所有 \(p \in P\)，\(T_{p}\) 是有限生成的挠 \(A_{p}\)-模，因而是有限长度的模（第四章，§2，第 5 节，命题 7 的推论 2）；将此长度记为 \(l_{p}(T)\)。现在，对于不包含 T 的湮灭子的所有 p，都有 \(T_{p} = 0\)，因而对几乎所有 p 都如此（§1，第 6 节，定理 4），这就说明下列定义是合理的：

定义 4. 若 T 是有限生成的挠 A-模，则除子：

\[
\chi (\mathbf {T}) = \sum_ {\mathfrak {p} \in \mathbf {P}} l _ {\mathfrak {p}} (\mathbf {T}). \mathfrak {p}.
\]

称为 T 的容度。

命题 10. (i) 设 \(0 \to T_1 \to T_2 \to T_3 \to 0\) 是有限生成挠 A-模的正合列。则

\[
\chi (\mathrm{T} _ {2}) = \chi (\mathrm{T} _ {1}) + \chi (\mathrm{T} _ {3}).
\]

\begin{enumerate}
\def\labelenumi{(\roman{enumi})}
\setcounter{enumi}{1}
\item
  若存在伪同构 \(f \colon \mathrm{T}_1 \to \mathrm{T}_2\)，则 \(\chi(\mathrm{T}_1) = \chi(\mathrm{T}_2)\)。
\item
  要有 \(\chi(T) = 0\)，充要条件是 \(T\) 为伪零模。
\end{enumerate}

根据定义 4，只需对每个 \(p \in P\) 考察所讨论挠模的 \(l_{p}\) 值。性质 (i) 由第二章，§2，第 4 节，定理 1 以及正合列中长度的可加性得出（代数，第二章，§1，第 10 节，命题 16），性质 (ii) 和 (iii) 则立即由第 4 节的定义得出。

推论。设 \(0 \to T_n \to T_{n-1} \to \cdots \to T_0 \to 0\) 是有限生成挠 A-模的正合列。则 \(\sum_{i=0}^{n} (-1)^i \chi(T_i) = 0\)。

根据第二章，§2，第 4 节，定理 1，这再次由 \(l_{\mathfrak{p}}\) 的相应性质得出（代数，第二章，§1，第 10 节，命题 16 的推论 3）。

回忆（第二章，§ 5，第 4 节），我们可以相对于同构关系谈论有限生成 A-模的类组成的集合 \(\mathbf{F}(\mathbf{A})\)；对于每个有限生成 A-模 M，令 \(\operatorname{cl}(\mathbf{M})\) 表示 \(\mathbf{F}(\mathbf{A})\) 中相应的元素；记 \(\mathbf{T}(\mathbf{A})\) 为 \(\mathbf{F}(\mathbf{A})\) 的子集，它由有限生成挠 A-模的类组成。显然，\(\chi\) 定义了从 \(\mathbf{T}(\mathbf{A})\) 到 \(\mathbf{D}(\mathbf{A})\) 的映射，仍记为 \(\chi\)，使得 \(\chi(\operatorname{cl}(\mathbf{T})) = \chi(\mathbf{T})\)。

命题 11. 设 G 是一个用加法记号书写的交换群，\(\phi: T(A) \to G\) 是一个映射；对于每个有限生成挠 A-模 T，我们也滥用语言写作 \(\phi(T) = \phi(\text{cl}(T))\)。假设满足下列条件：

\begin{enumerate}
\def\labelenumi{(\arabic{enumi})}
\item
  若 \(0 \to \mathrm{T}_1 \to \mathrm{T}_2 \to \mathrm{T}_3 \to 0\) 是有限生成挠 A-模的正合列，则 \(\phi(\mathrm{T}_2) = \phi(\mathrm{T}_1) + \phi(\mathrm{T}_3)\)。
\item
  若 \(\mathbf{T}\) 是伪零的，则 \(\phi (\mathbf{T}) = 0\)。
\end{enumerate}

则存在唯一同态 \(\theta: \mathbf{D}(\mathbf{A}) \to \mathbf{G}\)，使得 \(\phi = \theta \circ \chi\)。

由于 \(\chi(A/p) = p\) 对所有 \(p\) 成立，必有 \(\theta(p) = \phi(A/p)\) 对所有 \(p \in P\)，这证明了 \(\theta\) 的唯一性，因为 P 的元素构成 D(A) 的一组基。反之，设 \(P\)\(D(A)\)\(\theta\) 是从 D(A) 到 G 的同态，满足对所有 \(D(A)\)\(G\) 有 \(\theta(p) = \phi(A/p)\)，证明它满足要求。为此，写成\(p \in P\)

\[
\psi (\mathbf {T}) = \phi (\mathbf {T}) - \theta (\chi (\mathbf {T}))
\]

对每个有限生成挠 A-模 T；显然若以 \(\phi\) 替换为 \(\psi\)，条件 (1) 和 (2) 仍满足。另一方面，若 ，则 \(\psi(A/p) = 0\)；若 \(p \in P\)\(p\) 是一个 \(\neq 0\) 且不在 \(P\) 中的素理想，则 A/p 的湮灭子不包含于 \(A/p\)\(P\) 中任何理想，因而（第 4 节，定理 5）A/p 是伪零的，所以 \(A/p\)\(\psi(A/p) = 0\)。既然如此，每个有限生成挠 A-模 T 都有一个分解列，其因子同构于形如 \(A/p\) 的 A-模，其中 \(p \in \text{Supp}(T)\)（第四章，§ 1，第 4 节，定理 1 和定理 2），并且由于 T 是挠模，有 \(p \neq 0\)。对该分解列的长度作归纳，结合 \(T\)\(\psi\) 的性质 (1)，得到（根据定义）\(\psi(T) = 0\)。

* 我们可以像第 4 节那样，考虑由有限生成挠 A-模组成的范畴 \(\mathcal{T} / \mathcal{T}'\)\(\mathcal{T}\) 关于有限生成伪零挠 A-模组成的满子范畴 \(\mathcal{T}'\) 的商范畴 \(\mathcal{T} / \mathcal{T}'\)。用阿贝尔范畴的语言，命题 11 表明，阿贝尔范畴  的 Grothendieck 群典范同构于 D(A)。 *

命题 12. 若 a 是 A 的非零理想，\(\neq 0\)

\[
\chi (\mathrm{A} / \mathfrak {a}) = \chi ((\mathrm{A}: \mathfrak {a}) / \mathrm{A}) = \operatorname{div} \mathfrak {a}.
\]

令 \(p \in P\)。则 \(aA_p = p^{n_p}A_p\)，其中 \(n_p \geqslant 0\)，因为 \(A_p\) 是离散赋值环。由于 \((A/a)_p = A_p/aA_p\)，\(l_p(A/a) = n_p\)，从而 \(\chi(A/a) = \sum_{p \in P} n_p p = \text{div } a\)（\(\S 1, \text{ no. } 4, \text{ Proposition } 7\)）。

另一方面，\((\mathbf{A}:\mathfrak{a})_{\mathfrak{p}} = \mathbf{A}_{\mathfrak{p}}\colon \mathfrak{a}\mathbf{A}_{\mathfrak{p}} = \mathfrak{p}^{-n_{\mathfrak{p}}}\mathbf{A}_{\mathfrak{p}}\)，因此 \(l_{\mathfrak{p}}((\mathbf{A}:\mathfrak{a}) / \mathbf{A}) = n_{\mathfrak{p}}\)，同理可得结论。

